\section{Κατανεμημένοι Πίνακες Κατακερματισμού και Ευρετήρια RMI}
\label{sec:bg_dht_rmi}

\subsection*{Κλασικά DHT και Πρωτόκολλο Chord}

Οι Κατανεμημένοι Πίνακες Κατακερματισμού (Distributed Hash Tables --- DHTs) αποτελούν θεμελιώδη δομή για την κατανομή και αναζήτηση δεδομένων σε κατανεμημένα συστήματα χωρίς κεντρικό συντονιστή. Το πρωτόκολλο Chord~\cite{stoica2001chord} οργανώνει τους κόμβους και τα κλειδιά σε έναν \textit{λογικό δακτύλιο αναγνωριστικών} (identifier ring) μεγέθους $2^m$ bits, χρησιμοποιώντας κρυπτογραφική συνάρτηση κατακερματισμού (π.χ.\ SHA-1). Κάθε κόμβος $n$ διαχειρίζεται τα κλειδιά $k$ για τα οποία ισχύει $\text{hash}(k) \in (\text{pred}(n),\, n]$, σύμφωνα με την αρχή \textit{συνεπούς κατακερματισμού} (consistent hashing). Τυπικά, ο κόμβος υπεύθυνος για ένα κλειδί $k$ ορίζεται ως:
\begin{equation}
    \text{successor}(k) = \min\{n \in \mathcal{N} \mid n \geq k \pmod{2^m}\},
\end{equation}
όπου $\mathcal{N}$ είναι το σύνολο των ενεργών κόμβων. Κάθε κόμβος διατηρεί έναν \textit{πίνακα δακτύλων} (finger table) $F_n[i] = \text{successor}(n + 2^{i-1})$ για $i = 1, \ldots, m$, ο οποίος επιτρέπει δρομολόγηση ερωτήματος σε $O(\log N)$ βήματα, όπου $N$ ο αριθμός ενεργών κόμβων.

\paragraph{Εικονικοί κόμβοι (vnodes) και εξισορρόπηση φορτίου.}
Η χρήση \textit{εικονικών κόμβων} (virtual nodes, vnodes) αντιμετωπίζει την ανομοιόμορφη κατανομή κλειδιών που προκαλεί ο πεπερασμένος αριθμός φυσικών κόμβων~\cite{stoica2001chord}. Κάθε φυσικός κόμβος αντιστοιχίζεται σε $V$ εικονικές θέσεις στον δακτύλιο ($V \approx 100{-}300$ στην πράξη), με αποτέλεσμα η κατανομή του φορτίου να συγκλίνει προς ομοιομορφία με τυπική απόκλιση $O(1/\sqrt{VN})$. Ωστόσο, η αύξηση του $V$ επαυξάνει τη μνήμη των finger tables, δημιουργώντας ένα εμπορικό αντιστάθμισμα (trade-off) μεταξύ ισορροπίας φορτίου και κόστους συντήρησης μεταδεδομένων.

\subsection*{Από Κρυπτογραφικούς σε Μαθημένους Κατακερματιστές}

Αν και ο κλασικός συνεπής κατακερματισμός επιτυγχάνει ιδανική εξισορρόπηση φορτίου, καταστρέφει πλήρως τη χωρική γειτνίαση: δύο σημεία με απειροελάχιστη απόσταση στον φυσικό χώρο καταλήγουν σε εντελώς διαφορετικούς κόμβους, καθιστώντας αδύνατα τα ερωτήματα εύρους (range queries) και αναζήτησης κοντινότερων γειτόνων (k-NN). Για την υπέρβαση αυτού του περιορισμού, η πρόσφατη βιβλιογραφία εισήγαγε τα \textit{Ευρετήρια Μάθησης} (Learned Indexes).

\paragraph{Ιεραρχικό μοντέλο RMI.}
Οι Kraska et al.~\cite{kraska2018case} πρότειναν το \textit{Recursive Model Index} (RMI): ένα ιεραρχικό μοντέλο δύο σταδίων που αντικαθιστά το B-Tree. Στο πρώτο στάδιο, ένα καθολικό μοντέλο $f_0$ εκτιμά σε ποιον «κάδο» (bucket) ανήκει το κλειδί, ενώ στο δεύτερο στάδιο ένα τοπικό μοντέλο $f_1$ (αντιστοιχεί σε έναν εκ των $B$ ειδικευμένων μοντέλων) εκτιμά την ακριβή θέση:
\begin{equation}
    \hat{\text{pos}} = f_1^{\,(f_0(\text{key}))}\!(\text{key}),
    \qquad f_0 : \mathcal{K} \to \{0, 1, \ldots, B-1\}.
\end{equation}
Το RMI εκπαιδεύεται στατικά μία φορά πάνω στη γνωστή κατανομή κλειδιών, επιτυγχάνοντας σφάλμα πρόβλεψης εντός ενός μικρού σταθερού δελτίου $\delta$ με υψηλή πιθανότητα.

\paragraph{Πολυδιάστατα Learned Indexes και οι περιορισμοί τους.}
Η επέκταση σε πολλαπλές διαστάσεις αντιμετωπίζεται από το σύστημα FLOOD~\cite{nathan2020flood}, το οποίο εκμαθαίνει ανεξάρτητη διαμέριση ανά διάσταση και στη συνέχεια συνδυάζει τις διαμερίσεις σε ένα πλέγμα. Ωστόσο, το FLOOD αποτυγχάνει όταν η κατανομή των δεδομένων δεν είναι ευθυγραμμισμένη με τους άξονες (π.χ.\ διαγώνιες συσχετίσεις), καθώς το πλέγμα δεν μπορεί να αναπαραστήσει δομές εκτός αξόνων. Επίσης, το κόστος μνήμης κλιμακώνεται εκθετικά ως $O(B^d)$ με τον αριθμό διαστάσεων $d$.

\paragraph{LEAD DHT: Συνδυασμός Hilbert και RMI.}
Για να αντιμετωπίσει τους παραπάνω περιορισμούς, το σύστημα LEAD~\cite{wang2025lead} συνδυάζει την αναγωγή διαστάσεων μέσω καμπύλης Hilbert --- που διατηρεί τη χωρική γειτνίαση (βλ.\ Ενότητα~\ref{sec:bg_hilbert}) --- με ένα μονοδιάστατο RMI που εκτιμά τον υπεύθυνο κόμβο για κάθε κλειδί. Εικονικοί κόμβοι (vnodes) και ένας επιγραμμικός ελεγκτής PID συνεργάζονται για την ελαχιστοποίηση της ανισορροπίας φορτίου, επιτρέποντας αποδοτικά ερωτήματα εύρους και k-NN εντός του κατανεμημένου δακτυλίου.
